\section{Completez. Teorema de Bolzano-Weierstrass}

A lo largo de esta sección $(X,d)$ denotará un espacio métrico.
Comenzamos estableciendo unas propiedades básicas acerca de sucesiones en $\R$.

\begin{teorema}\label{pro-basi-suces}
Sean $\{x_n\}$ y $\{y_n\}$ sucesiones en $\R$ convergentes a $x$ y $y$ respectivamente. Entonces
\begin{enumerate}
    \item $\{x_n+y_n\}$ converge a $x+y$.
    \item Si $k\in \R$, entonces $\{k x_n\}$ converge a $kx$.
    \item Si existe un $M$ tal que $x_n \leq y_n$ para todo $n\geq M$, entonces $x\leq y$.  
        \item $\{x_ny_n\}$ converge a $xy$.
\end{enumerate}
\end{teorema}
\begin{proof}
\begin{enumerate}
    \item Este resultado se sigue de la definición, teniendo en cuenta que para todo $n$
    \[|(x_n+y_n)-(x+y)|\leq|x_n-x|+|y_n-y|.\]
    \item Sea $k\in \R$. Si $k=0$, entonces $\{k x_n\}$ es la sucesión constante $0$, la cual claramente converge a $0=kx$. Por otro lado, si $k\neq 0$. Sea $\varepsilon >0$. Como $\frac{\varepsilon}{|k|}>0$, existe un $N$ tal que $|x_n-x|<\frac{\varepsilon}{|k|}$ para todo $n\geq N$. Por tanto, si $n\geq N$, tenemos
    \[|kx_n-kx|=|k||x_n-x|\leq |k|\frac{\varepsilon}{|k|}=\varepsilon.\]
    \item Supongamos por el contrario que $y<x$. Sea $\varepsilon=\frac{x-y}{2}>0.$ Luego existe un $n>M$ tal que $|x_n-x|<\varepsilon$ y  $|y_n-y|<\varepsilon$. De donde,
    $x-x_n<\varepsilon$ y $y_n-y<\varepsilon$. Luego $\frac{x+y}{2}=x-\varepsilon<x_n$ y $y_n<\varepsilon+y=\frac{x+y}{2}$. Por tanto $y_n<x_n$, para algún $n>M$. Absurdo.
    \item Como $\{x_n\}$ es convergente, entonces es acotada. Luego existe un $C$ tal que $|x_n|\leq C$ para todo $n\in  \N.$
    \begin{align*}
    |x_ny_n-xy|& =|x_ny_n-x_ny-xy+x_ny| \\
    & = |x_n(y_n-y)+y(x_n-x)| \\
    & \leq |x_n(y_n-y)|+|y(x_n-x)|\\
    &=|x_n||y_n-y|+|y||x_n-x|\\
    &\leq C|y_n-y|+|y||x_n-x|.
    \end{align*}
        De los pasos anteriores $\lim_{n\to \infty}C|y_n-y|+|y||x_n-x|=0 .$ Por tanto $\lim_{n\to \infty}|x_ny_n-xy|=0.$
    Es decir $x_ny_n\to xy.$
\end{enumerate}
\end{proof}

\begin{definicion}
Dada una sucesión $\{x_n\}$ en $X$, decimos que es de Cauchy, si para todo $\varepsilon >0$, existe un natural $N=N(\varepsilon)$ tal que \[d(x_n, x_m) <\varepsilon, \qquad \text{para todo $n,m \geq N$}.\]
\end{definicion}
Hablando en lenguaje ordinario, podría decirse que una sucesión de Cauchy es una sucesión que quiere ser convergente, pues los valores de ésta son cada vez más cercanos unos de otros. Sin embargo esta condición (de ser de Cauchy) no implica convergencia.

\begin{example}
Considere la sucesión $\left\{ \frac{1}{n}\right\}$, la cual es de Cauchy en $X=(0,1]$ pero no es convergente en $X$.
\end{example}

\begin{teorema}
Toda sucesión convergente en $X$ es de Cauchy.
\end{teorema}
\begin{proof}
Sea $\left\{x_n \right\}$ una sucesión en $X$ que converge a $x\in X$. Teniendo en cuenta la desigualdad
\[d(x_n, x_m)\leq d(x_n, x)+d(x, x_m), \quad \text{para todo $m,n\in \N$}, \]
se sigue el resultado.
\end{proof}

\begin{definicion}
Sea $A\subseteq X$.
\begin{enumerate}
    \item La clausura de $A$ es $\overline{A} :=A\cup A'$.
    \item Decimos que $A$ es denso en $X$ si $\overline{A}=X$.
    \item Sea $(Y, d_Y)$ un espacio métrico  y $\Phi: X\to Y$ una biyección. Decimos que $\Phi $ es una isometría de $X$ en $Y$ si
    \[d(x_1,\, x_2)=d_Y(\Phi (x_1),\, \Phi (x_2)), \qquad \text{para todo $x_1,x_2 \in X$}.\]
    \item Si existe una isometría de $X$ en $Y$, decimos que $X$ es isométrico a $Y$ ($X\cong Y$).
    \item Decimos que $X$ es completo si toda sucesión de Cauchy en $X$, es convergente en $X$.
\end{enumerate}
\end{definicion}

\begin{teorema}\label{cauchy-acota}
Toda sucesión de Cauchy en $\R$ es acotada.
\end{teorema}
\begin{proof}
Sea $\left\{x_n \right\}$ una sucesión de Cauchy en $X$. Por tanto existe un $N$ tal que para todo $m\geq N$, $x_m\in B(x_N, 1).$ Sea $R=2\max\{1, \, d(x_1, x_N),\dots, d(x_{N-1}, x_N)\}.$ Con esto tenemos que $x_n \in B(x_N, R)$ para todo $n\in \N.$
\end{proof}

\begin{remark}
En lo que sigue usaremos un hecho fundamental de los números reales, que se toma como axioma en la construcción de $\R$: \emph{la propiedad del supremo}. Esto es, todo subconjunto no vacío de $\R$ que es acotado superiormente posee supremo en $\R$. (Este hecho ya se usó, por ejemplo, al definir la métrica del supremo en $B(\N)$.)
\end{remark}

\begin{teorema}[Bolzano-Weierstrass]\label{bolzano-weierstrass}
Toda sucesión acotada de números reales posee una subsucesión convergente.
\end{teorema}
\begin{proof}
Sea $\{x_n\}$ una sucesión contenida en un intervalo $[A,B]$. Construimos por recursión intervalos encajados $I_k=[a_k,b_k]$ tales que, para cada $k$, el conjunto de índices $\{n\in\N : x_n\in I_k\}$ es infinito y la longitud de $I_k$ es $(B-A)/2^k$.

Tomamos $I_0=[A,B]$. Supongamos construido $I_{k-1}=[a,b]$ que contiene a $x_n$ para infinitos índices $n$, y sea $m=(a+b)/2$ su punto medio. Como $I_{k-1}=[a,m]\cup[m,b]$ y contiene a $x_n$ para infinitos índices, al menos una de las dos mitades contiene a $x_n$ para infinitos índices; definimos $I_k$ como una de tales mitades.

Escogemos ahora, recursivamente, índices $n_1<n_2<\cdots$ tales que $x_{n_k}\in I_k$: dado $n_{k-1}$, ello es posible porque $I_k$ contiene a $x_n$ para infinitos índices $n$.

La sucesión $\{a_k\}$ es creciente y acotada superiormente (por $B$); por la propiedad del supremo, $x:=\sup\{a_k : k\in\N\}$ existe en $\R$. Afirmamos que $a_k\le x\le b_k$ para todo $k$. La primera desigualdad es la definición de supremo. Para la segunda, observemos que cada $b_k$ cota superiormente a $\{a_j : j\in\N\}$: si $j\le k$ entonces $a_j\le a_k\le b_k$, y si $j\ge k$ entonces $a_j\le b_j\le b_k$. Por tanto $x\le b_k$.

Finalmente, como $x_{n_k}\in I_k=[a_k,b_k]$ y $a_k\le x\le b_k$,
\[
|x_{n_k}-x|\ \le\ b_k-a_k\ =\ \frac{B-A}{2^k}\ \longrightarrow\ 0,
\]
de modo que $x_{n_k}\to x$. Por lo tanto, $\{x_n\}$ posee una subsucesión convergente.
\end{proof}

\begin{corolario}\label{compacto-intervalo}
Para todos $a<b$, el intervalo $[a,b]$ es un subconjunto compacto de $\R$ (con la métrica usual).
\end{corolario}
\begin{proof}
Sea $\{x_n\}$ una sucesión en $[a,b]$. En particular, $\{x_n\}$ es acotada y, por el Teorema de Bolzano-Weierstrass, posee una subsucesión $\{x_{n_k}\}$ convergente, con límite $x\in\R$. Como $[a,b]$ es cerrado y $x_{n_k}\in[a,b]$ para todo $k$, la caracterización secuencial de los conjuntos cerrados da $x\in[a,b]$. Por lo tanto, toda sucesión en $[a,b]$ posee una subsucesión convergente con límite en $[a,b]$; es decir, $[a,b]$ es compacto.
\end{proof}

\begin{teorema}\label{R-completo}
$\R$ es un espacio métrico completo.
\end{teorema}
\begin{proof}
Sea $\{x_n \}_n$ una sucesión de Cauchy en $\R$. Por Teorema \ref{cauchy-acota}, existe un $R>0$ tal que $x_n \in [-R, \, R]$, para todo $n$. Como $[-R, \, R]$ es compacto (Corolario \ref{compacto-intervalo}), esta sucesión admite una  subsucesión convergente $\{x_{n_k}\}_k$. Supongamos que $x$ es su límite. Sea $\varepsilon >0$. Entonces existen  $N_1$ y $N_2$ tales que
\[ | x_n- x_m|<\frac{\varepsilon}{2} \quad \text{para todo $m,n\geq N_1$}\]
y
\[ | x_{n_k}- x|<\frac{\varepsilon}{2} \quad \text{para todo $k\geq N_2$}.\]
Fijemos un $k\geq \max\{N_1, N_2\}$. Entonces, $n_k\geq N_1$ y $n_k\geq N_2$. Así que teniendo en cuenta las desigualdades anteriores
\[ | x_n- x | \leq | x_n- x_{n_k} |+| x_{n_k}- x | <\frac{\varepsilon}{2}+\frac{\varepsilon}{2}={\varepsilon}, \quad \text{para todo $n \geq N_1$}.\]
Es decir, $x_n \to x$ cuando $n\to \infty .$
\end{proof}

\begin{definicion}
Decimos que $X$ es totalmente acotado si para todo $\varepsilon >0$, existe un número finito de puntos de $X$, $x_1, \dots, x_k$ tal que $X=\bigcup _{i=1}^k B(x_i, \varepsilon).$
\end{definicion}

\begin{teorema}
$X$ es compacto si y sólo si $X$ es completo y totalmente acotado.
\end{teorema}
